\cvsection{Research}

\cvsubsection{Grants and Contracts}
\begin{cvenum}
  \item National Ocean Partnership Program via Woods Hole Oceanographic Institution, ``EPIC-DAUG: An Enhanced Propulsion Integrated Capability: Deep Autonomous Underwater Glider for Shore Launch/Recovery and High-Endurance Unattended Mission.'' PI: B.\ Bingham. \$78,705 (8/1/2020--9/30/2021).
  \item Office of Naval Research, ``Consortium for Robotics and Unmanned Systems Education and Research (CRUSER).'' PI: B.\ Bingham. \$4,000,000 (1/1/2020--12/31/2020).
  \item Defense Advanced Research Projects Agency, ``Virtual Environments to Support Development of Underwater Autonomy.'' PI: B.\ Bingham. \$280,133 (3/1/2020--12/31/2020).
  \item Office of Naval Research, ``Virtual Environment to Support Ocean Robotics.'' PI: B.\ Bingham. \$150,718 (8/1/2020--12/31/2020).
  \item Office of Naval Research, ``Consortium for Robotics and Unmanned Systems Education and Research (CRUSER).'' PI: B.\ Bingham. \$4,000,000 (10/1/2018--12/31/2019).
  \item Office of Naval Research, ``Virtual Environment to Support Ocean Robotics.'' PI: B.\ Bingham. \$117,870 (7/1/2017--12/31/2019).
  \item Space and Naval Warfare Systems Center Pacific, ``Development of Short-Course Materials for Robotics.'' PI: B.\ Bingham. \$32,685 (10/1/2016--9/30/2017).
  \item Helmsley Charitable Trust, ``The Vertically Integrated Projects (VIP) Consortium: Achieving Systematic Reform of Higher Education in STEM Fields.'' PI: Aaron Ohta (UHM); Co-PIs: B.\ Bingham, Z.\ Trimble, D.\ Garmire, M.\ Nejhad and Wayne Shiroma. \$270,000 (1/1/2015--12/31/2017).
  \item Army Communications Electronics Command, ``Hawaii Undersea Military Munitions Assessment: Advanced Data Processing.'' PI: M.\ Edwards (UHM); Co-PI: B.\ Bingham. \$469,889 (9/10/2014--9/26/2015).
  \item Makai Ocean Engineering, ``Array Burial Vehicle: Phase II, Option 2.'' PI: B.\ Bingham. \$204,735 (6/1/2014--2/1/2015).
  \item National Science Foundation (NSF), ``Autonomous, Low Cost Platform for Seafloor Geodesy.'' PI: James Foster (UHM); Co-PI: B.\ Bingham. \$342,307 (8/15/2013--7/31/2014).
  \item National Science Foundation (NSF), ``Collaborative Research: Distributed Heterogeneous Ocean Robots for Detecting and Monitoring Oil Plumes.'' PI: B.\ Bingham (collaboration with Stevens Institute of Technology, Dr.\ Yi Guo). \$199,336 (9/1/2012--8/31/2015).
  \item National Oceanographic and Atmospheric Administration (NOAA) through University Corporation for Atmospheric Research (UCAR), ``NOAA OER ROV Software Development.'' PI: B.\ Bingham. \$35,108 (10/1/2012--9/30/2013).
  \item Makai Ocean Engineering, ``Array Burial Vehicle: Phase II.'' PI: B.\ Bingham. \$151,444 (8/1/2012--8/31/2014).
  \item Makai Ocean Engineering, ``Array Burial Vehicle: Phase I Option.'' PI: B.\ Bingham. \$20,591 (4/27/2012--9/30/2012).
  \item Department of Homeland Security, ``Unmanned Port Security Autonomous System for Maritime and Port Security.'' PI: B.\ Bingham; Co-PIs: H.\ Singh (WHOI) and R.\ Camilli (WHOI). \$499,968 (2/12/2010--2/11/2013).
  \item Pukoa Scientific LLC, ``Navigation and Sensor Fusion for Image Tracking from a UAV.'' PI: B.\ Bingham. \$29,585 (5/1/2011--6/30/2012).
  \item Makai Ocean Engineering, ``Underwater Sensor System for Autonomous Burial and Operation.'' PI: B.\ Bingham. \$24,534 (7/11/2011--1/27/2012).
  \item National Oceanographic and Atmospheric Administration (NOAA) through University Corporation for Atmospheric Research (UCAR), ``Okeanos ROV Research Cruise.'' PI: B.\ Bingham. \$64,497 (4/11/2011--8/31/2011).
  \item National Defense Center of Excellence for Research in Ocean Sciences (CEROS), ``Autonomous Gateway Wave Glider Supporting Acoustic Tracking and Communication for a Portable Underwater Tracking Range.'' PI: B.\ Bingham. \$268,282 (8/15/2010--12/31/2011).
  \item National Oceanographic and Atmospheric Administration (NOAA) through University Corporation for Atmospheric Research (UCAR), ``R/V Okeanos ROV Research Cruise.'' PI: B.\ Bingham. \$18,511 (7/12/2010--8/11/2010).
  \item Office of Naval Research (ONR), 2008 Defense University Research Instrumentation Program (DURIP), ``Advanced Ultra-Short Baseline Navigation Systems for the National Deep Submergence Facility: Advancing Expeditionary Science and Rapid Event Response in Ocean Science and Engineering.'' PI: A.\ Bowen; Co-PIs: B.\ Bingham and L.\ Whitcomb (JHU). \$637,900 (4/15/2008--4/14/2009).
  \item Schlumberger, ``SCOPE Senior Design Project: Autonomous Surface Craft.'' PI: B.\ Bingham. \$50,000 (9/1/2007--6/1/2008).
  \item U.S.\ Coast Guard Research and Development Center, ``Rapid Detection of Submerged Heavy Oil Spills with Real-Time Mass Spectrometry, Fluorometry, and Acoustic Navigation.'' PI: R.\ Camilli; Co-PI: B.\ Bingham.
  \item NOAA Office of Ocean Exploration, ``Autonomous Rapid High Resolution Mapping of Ancient Deep Water Shipwrecks and Geologic Features.'' PI: B.\ Foley (WHOI); Co-PIs: H.\ Singh, R.\ Camilli, D.\ Mindell, B.\ Bingham, Ryan Eustice, Katerina Delaporta and Dimitris Sakellariou.
  \item Woods Hole Oceanographic Institution --- Tropical Research Initiative, ``Rapid Swath Mapping of Reef Ecology and Associated Water Column Chemistry in the Gulf of Chiriquí, Panamá.'' PI: R.\ Camilli; Co-PIs: L.\ Camilli, O.\ Pizzaro and B.\ Bingham.
  \item NSF --- Marine Geology and Geophysics, ``SST: Smarter Sensing: Integrating Payloads, Navigation, and Vehicles for Underwater Exploration.'' PI: H.\ Singh; Co-PIs: B.\ Bingham, J.\ Whelan and R.\ Camilli. \$763,134 (9/1/2004--8/31/2007).
\end{cvenum}

\cvsubsection{Journal Publications}
\printbibliography[keyword=journal, heading=none, resetnumbers=true]

\cvsubsection{Non-Technical Publications}
\printbibliography[keyword=nontechnical, heading=none, resetnumbers=true]

\cvsubsection{Conference Publications}
\printbibliography[keyword=conference, heading=none, resetnumbers=true]

\cvsubsection{Conference Presentations and Abstracts}
\printbibliography[keyword=presentation, heading=none, resetnumbers=true]

\cvsubsection{Book Chapters}
\printbibliography[keyword=bookchapter, heading=none, resetnumbers=true]

\cvsubsection{Expeditions and Field Experience}
\begin{cvenum}
  \item Hydrothermal Dynamics of Yellowstone Lake, Summers 2016--2018. National Science Foundation via Global Foundation for Ocean Exploration.
  \item Hohonu Moana: Exploring the Deep Waters Off Hawaii, August 2015. NOAA Ocean Exploration and Research, ROV Team Lead.
  \item Océano Profundo 2015: Exploring Puerto Rico's Seamounts, Trenches, and Troughs, April 2015. NOAA Ocean Exploration and Research, ROV Team Lead.
  \item Exploring Atlantic Canyons and Seamounts, April 2014. ROV Team Lead.
  \item Northeast U.S.\ Canyons Expedition and Shakedown Cruise, May and July 2013. NOAA Ocean Exploration and Research, ROV Team Lead.
  \item Sentry AUV Exploration and Testing, Blake Ridge Diapir, July 2012. National Science Foundation and NOAA Ocean Exploration and Research. Chief Scientist: Cindy Van Dover.
  \item GALREX 2011: Galápagos Rift Expedition, Summer 2011. NOAA Ocean Exploration and Research.
  \item INDEX 2010: Indonesia--USA Deep-Sea Exploration of the Snagihe Talaud Region, Summer 2010. NOAA Ocean Exploration and Research.
  \item Deep Chemosynthetic Community Characterization: ROV survey and sampling, June 2007. Woods Hole Oceanographic Institution, Penn State, TDI-Brooks, University of Georgia, Texas A\&M.
  \item Survey of Classical Greek Shipwrecks, ancient shipwrecks in the Aegean Sea, July 2006. Dr.\ Brendan Foley, Woods Hole Oceanographic Institution.
  \item Chios Shipwreck Survey: Autonomous underwater vehicle surveys of ancient shipwrecks in the Aegean Sea, July 2005. Dr.\ Hanumant Singh and Dr.\ Brendan Foley, Woods Hole Oceanographic Institution.
  \item Ridge 2000: ROV exploration of Lau Basin hydrothermal vent sites, June 2005. Dr.\ Chuck Fisher, Penn State University.
  \item Kane Fault: ROV sampling of Kane fault section of the mid-Atlantic ridge, December 2004. Dr.\ Maurice Tivey, Woods Hole Oceanographic Institution.
  \item STAG: ROV bathymetry survey of Trans-Atlantic Geotraverse (TAG) hydrothermal mound, mid-Atlantic ridge, November 2004. Dr.\ Rob Reves-Sohn, Woods Hole Oceanographic Institution.
  \item SeaBreeze: Physical oceanography and mapping surveys of the Endeavor hydrothermal vent area using the autonomous benthic explorer (ABE), July 2004. Dr.\ Russell McDuff, University of Washington.
  \item Mountains in the Sea: Biological exploration and survey of the New England Seamounts using the ROV Hercules, May 2004. Dr.\ Les Watling, University of Maine, URI, IFE, WHOI.
  \item Leon Levy Expedition: Roman shipwreck survey, excavation, and recovery at Skerki Bank, Mediterranean Sea, Summer 2003. Dr.\ Robert Ballard, URI, IFE, WHOI, and MIT.
  \item Black Sea Expedition: Remotely operated vehicle survey and excavation of ancient shipwrecks off the coast of Turkey, Summer 2003. Dr.\ Robert Ballard, URI, IFE, WHOI, and MIT.
  \item USS Monitor: AUV survey of the USS Monitor off the coast of Cape Hatteras, NC, May 2003. MIT-DeepArch, WHOI, and NOAA Monitor Marine Sanctuary.
  \item USS Monitor: Sub-bottom sonar survey of the USS Monitor's gun turret, Cape Hatteras, NC, August 2001. NOAA Monitor Marine Sanctuary and MIT-DeepArch.
  \item AUV Archaeology: Archaeological survey of Greek shipwreck using autonomous underwater vehicles, Nisyros, Greece, July 2001. Bluefin Robotics, MIT-DeepArch and the MIT AUVlab.
  \item Defense: High-frequency sub-bottom profiling sonar survey of Revolutionary War privateer, October 2001. MIT-DeepArch, Plymouth State College, and U-Maine Darling Marine Center.
  \item CSS Hunley: Sub-bottom profiling of Civil War submarine in on-site preservation, Charleston, South Carolina, June 2000. MIT-DeepArch, Marine Sonics Technology, and the Friends of the Hunley.
\end{cvenum}
